\documentclass[aspectratio=169, 11pt]{beamer}
\usetheme{metropolis}

\usepackage{graphicx}
\usepackage{booktabs}
\usepackage{xcolor}
\usepackage{appendixnumberbeamer}
\usepackage{listings}
\usepackage{qrcode}

% Custom colors
\definecolor{alertred}{RGB}{204, 36, 29}
\definecolor{okgreen}{RGB}{40, 140, 40}
\definecolor{noteblue}{RGB}{30, 90, 160}
\definecolor{codebg}{RGB}{242, 242, 242}

\setbeamercolor{alerted text}{fg=alertred}

\graphicspath{{figs/}}

% Code listing style
\lstset{
  basicstyle=\ttfamily\footnotesize,
  backgroundcolor=\color{codebg},
  frame=none,
  breaklines=true,
  showstringspaces=false,
  columns=fullflexible,
  keepspaces=true,
  xleftmargin=0.5em,
  xrightmargin=0.5em,
  aboveskip=0.4em,
  belowskip=0.2em,
}

\title{Linux \& HPC Basics}
\subtitle{Getting Up and Running on the Cluster}
\author{Caleb Sy}
\date{March 2026}
\institute{PhD Research Briefing}

\begin{document}

\maketitle

%% ============================================================
%% PART I: THE LINUX KERNEL
%% ============================================================

\section{The Linux Kernel}

% ---- Slide 2: What Is Linux? ----
\begin{frame}{What Is Linux?}
  \begin{columns}[T]
    \column{0.52\textwidth}
      The \textbf{kernel} is the core of the operating system ---
      the layer between hardware and your programs.
      \vspace{0.4em}
      \begin{itemize}
        \item Manages CPU scheduling, memory, storage, networking, and I/O
        \item Every process running on the cluster goes through it
        \item You interact with it through a \textbf{shell} (bash)
      \end{itemize}
      \vspace{0.4em}
      \begin{block}{The software stack}
        \textbf{Hardware} $\to$ \textbf{Linux Kernel} $\to$ \textbf{Bash Shell} $\to$ \textbf{Your Code}
      \end{block}

    \column{0.44\textwidth}
      \begin{block}{On Henry2}
        When you SSH into the cluster, you land in a bash shell.
        Commands you type go to the kernel, which runs them on the hardware.
      \end{block}
      \vspace{0.5em}
      \begin{itemize}
        \item \textbf{Linux} = the kernel + surrounding tools
        \item \textbf{Bash} = the shell you type commands into
        \item \textbf{LSF} = the job scheduler that sits on top
        \item \textbf{conda} = environment manager
      \end{itemize}
  \end{columns}
\end{frame}

% ---- Slide 3: Why Linux for HPC ----
\begin{frame}{Why Linux for HPC}
  \begin{itemize}
    \item \textbf{Free and open source} --- every HPC cluster in the world runs Linux
    \item \textbf{Stable under heavy load} --- built for long-running, multi-user environments
    \item \textbf{Shell scripting is built in} --- automate repetitive tasks with a text file
    \item \textbf{Environment management} --- \texttt{conda activate}, \texttt{module load} work natively
    \item \textbf{Remote access} --- SSH from any machine, no GUI required
  \end{itemize}
  \vspace{0.5em}
  \begin{alertblock}{Windows paths do not work here}
    \texttt{C:\textbackslash{}Users\textbackslash{}calisy\textbackslash{}...} is a Windows path.
    On Linux everything uses forward slashes and starts from \texttt{/}.\\
    Your home: \texttt{\$HOME} \quad Scratch storage: \texttt{/share/\$GROUP/\$USER/}
  \end{alertblock}
\end{frame}

%% ============================================================
%% PART II: BASIC COMMANDS
%% ============================================================

\section{Basic Commands}

% ---- Slide 4: Commands Reference ----
\begin{frame}{Commands You Will Use Every Day}
  \begin{table}
    \centering
    \small
    \begin{tabular}{lll}
      \toprule
      \textbf{Command} & \textbf{What it does} & \textbf{Example} \\
      \midrule
      \texttt{ls}   & List files in a directory                    & \texttt{ls -lh} \\
      \texttt{cd}   & Change directory                              & \texttt{cd scripts/} \\
      \texttt{cp}   & Copy a file or directory                      & \texttt{cp a.py b.py} \\
      \texttt{mv}   & Move or rename a file                         & \texttt{mv old.py new.py} \\
      \texttt{echo} & Print text or a variable to the terminal      & \texttt{echo \$MY\_VAR} \\
      \texttt{cat}  & Print a file's full contents to screen        & \texttt{cat out.txt} \\
      \texttt{less} & Scroll through a file (press \texttt{q} to quit) & \texttt{less log.txt} \\
      \texttt{nano} & Simple in-terminal text editor                & \texttt{nano script.sh} \\
      \bottomrule
    \end{tabular}
  \end{table}
  \vspace{0.4em}
  {\small Append \texttt{--help} to any command for usage info: \quad \texttt{ls --help} \quad \texttt{cp --help}}
\end{frame}

%% ============================================================
%% PART III: BASH SCRIPTING
%% ============================================================

\section{Bash Scripting}

% ---- Slide 5: Anatomy of a Bash Script ----
\begin{frame}[fragile]{Anatomy of a Bash Script}
  \begin{columns}[T]
    \column{0.54\textwidth}
      \begin{lstlisting}[language=bash]
#!/bin/bash
# this is a comment

# source a config file (imports its variables)
. ./experiment_grid_params.sh

# variable assignment (no spaces around =)
MY_VAR="hello world"

# print a variable
echo $MY_VAR

# loop over an array
for val in "${my_array[@]}"; do
    echo "Running: $val"
    python run_model.py $val
done
      \end{lstlisting}

    \column{0.42\textwidth}
      \begin{block}{Key rules}
        \begin{itemize}
          \item \texttt{\#!/bin/bash} must be line~1
          \item No spaces around \texttt{=} in assignments
          \item Use \texttt{\$VAR} or \texttt{\$\{VAR\}} to expand a variable
          \item Make executable before first run:\\
            \texttt{chmod +x script.sh}
        \end{itemize}
      \end{block}
      \vspace{0.4em}
      \begin{block}{Three ways to run}
        \texttt{./script.sh} \hfill subprocess\\
        \texttt{bash script.sh} \hfill same\\
        \texttt{. ./script.sh} \hfill \textit{source} (imports variables into current shell)
      \end{block}
  \end{columns}
\end{frame}

% ---- Slide 6: File Naming Conventions ----
\begin{frame}{File Naming Conventions}
  \begin{table}
    \centering
    \small
    \begin{tabular}{lll}
      \toprule
      \textbf{Pattern} & \textbf{Example} & \textbf{Purpose} \\
      \midrule
      \texttt{verb\_noun.sh}          & \texttt{run\_experiment.sh}            & Action scripts --- run these to do things \\
      \texttt{check\_*.sh}            & \texttt{check\_job\_status.sh}         & Monitoring and utility scripts \\
      \texttt{*\_params.sh}           & \texttt{experiment\_grid\_params.sh}   & Config files sourced by other scripts \\
      \texttt{*\_log\_MMDD\_HHMM.txt} & \texttt{failed\_results\_gather\_log\_0226\_1351.txt} & Timestamped error logs \\
      \bottomrule
    \end{tabular}
  \end{table}
  \vspace{0.5em}
  \begin{columns}[T]
    \column{0.48\textwidth}
      \begin{alertblock}{Never use spaces in filenames}
        \texttt{my file.sh} breaks shell scripts.\\
        \texttt{my\_file.sh} always works.\\
        Use underscores everywhere on HPC.
      \end{alertblock}

    \column{0.48\textwidth}
      \begin{block}{Case convention}
        Files and directories: \texttt{lowercase\_with\_underscores}\\
        Environment variables and constants: \texttt{UPPERCASE}
      \end{block}
  \end{columns}
\end{frame}

%% ============================================================
%% PART IV: RUNNING PYTHON
%% ============================================================

\section{Running Python}

% ---- Slide 7: Calling Python ----
\begin{frame}[fragile]{Calling Python from the Command Line and Bash}
  \begin{columns}[T]
    \column{0.47\textwidth}
      \textbf{From the command line:}
      \begin{lstlisting}[language=bash]
# basic run
python script.py

# with positional arguments
python script.py arg1 arg2

# as used in the model
python v13_calibration_wrapper.py \
    372 0.0
      \end{lstlisting}
      \vspace{0.3em}
      {\small Read args in the script with\\
      \texttt{import sys}\\
      \texttt{sys.argv[1]}, \texttt{sys.argv[2]}, \ldots}

    \column{0.49\textwidth}
      \textbf{From inside a bash script:}
      \begin{lstlisting}[language=bash]
adjustment=0.98
qps_shift=0.1
storage_shift=0.0

python v13_calibration_wrapper.py \
    ${adjustment} \
    ${qps_shift} \
    ${storage_shift}
      \end{lstlisting}
      \vspace{0.3em}
      {\small \texttt{\$\{...\}} expands to the variable's value at
      run time --- this is how the calibration grid iterates over
      all parameter combinations.}
  \end{columns}
  \vspace{0.3em}
  \begin{block}{This is exactly what \texttt{run\_experiment.sh} does}
    It loops over all \texttt{(adjustment, qps\_shift, storage\_shift)} combinations
    and calls \texttt{bsub "python v13\_calibration\_wrapper.py \$\{...\}"} for each one.
  \end{block}
\end{frame}

%% ============================================================
%% PART V: HPC WITH LSF
%% ============================================================

\section{HPC with LSF}

% ---- Slide 8: LSF Submission Script ----
\begin{frame}[fragile]{LSF Job Submission with \texttt{bsub}}
  \begin{lstlisting}[language=bash]
bsub -W 30 -n 1 -o "out.%J" -e "err.%J" "python v13_calibration_wrapper.py 372 0.0"
  \end{lstlisting}
  \vspace{0.2em}
  \begin{table}
    \centering
    \small
    \begin{tabular}{lll}
      \toprule
      \textbf{Flag} & \textbf{Meaning} & \textbf{Typical value} \\
      \midrule
      \texttt{-W 30}      & Wall time limit (minutes)       & \texttt{30} for model run; \texttt{5} for data setup \\
      \texttt{-n 1}       & Number of CPU cores             & \texttt{1} (model is single-threaded) \\
      \texttt{-o out.\%J} & Write stdout to this file       & \texttt{\%J} expands to the LSF job ID \\
      \texttt{-e err.\%J} & Write stderr to this file       & Read this first when a job fails \\
      \bottomrule
    \end{tabular}
  \end{table}
  \vspace{0.3em}
  \begin{columns}[T]
    \column{0.48\textwidth}
      \textbf{Monitor your jobs:}
      \begin{itemize}
        \item \texttt{bjobs -r} \quad\quad running jobs
        \item \texttt{bjobs -d} \quad\quad recently finished
        \item \texttt{bhist -l \$JOBID} \quad full history
      \end{itemize}
    \column{0.48\textwidth}
      \begin{block}{If a job fails}
        Check \texttt{err.\$JOBID} first.
        Common causes: wall time exceeded (\texttt{-W} too short),
        wrong conda env, or a missing input file.
      \end{block}
  \end{columns}
\end{frame}

% ---- Slide 9: Calibration File Structure ----
\begin{frame}[fragile]{Calibration Experiment File Structure}
  \begin{columns}[T]
    \column{0.46\textwidth}
      \textbf{Script directory} (in the repo):
      \begin{lstlisting}
v13_run_scripts/
  experiment_grid_params.sh   <- 1. edit params
  folder_setup.sh             <- 2. create dirs
  run_data_setup.sh           <- 3. build .dat files
  run_experiment.sh           <- 4. submit model jobs
  gather_run_results.sh       <- 5. collect outputs
  copy_to_research_storage.sh <- 6. archive
  check_job_status.sh            (job monitor utility)
      \end{lstlisting}

    \column{0.50\textwidth}
      \textbf{Run directories on scratch} (created by step~2):
      \begin{lstlisting}
/share/$GROUP/$USER/
  prod_experiment/
    set_of_15_prod_adjustments/
      prodshift_0.970_qpsshift_0.0_.../
        model_data_inputs_5-2-25/
        pickles/
        v13_calibration_wrapper.py
        v13_run_datasetup_...py
        ...
      prodshift_0.971_qpsshift_0.0_.../
      ...  (one directory per parameter combo)
      \end{lstlisting}
  \end{columns}
  \vspace{0.3em}
  \begin{block}{Run scripts 1 through 6 in order}
    Each script sources \texttt{experiment\_grid\_params.sh} first,
    then loops over all parameter combinations and submits \texttt{bsub} jobs.
  \end{block}
\end{frame}

%% ============================================================
%% PART VI: RESOURCES
%% ============================================================

\section{Resources}

% ---- Slide 10: QR Code ----
\begin{frame}{Get Up and Running: YouTube Series}
  \begin{center}
    \vspace{0.2em}
    \qrcode[height=4.2cm]{https://www.youtube.com/watch?v=YwqpWSOWHNM\&list=PLskHZ4tWojE1NZKZiF_XRVztGAJatUKFt\&index=2}

    \vspace{0.5em}
    {\large \textbf{Linux Command Line for Beginners}}

    \vspace{0.3em}
    {\small Full playlist --- start at video~2 for the fastest path to being productive}

    \vspace{0.3em}
    {\footnotesize \texttt{https://www.youtube.com/watch?v=YwqpWSOWHNM}}
  \end{center}
\end{frame}

\end{document}
